\documentclass{ximera}

\input{../../../preamble.tex}


\definecolor{fvdnavy}{HTML}{071D43}
\definecolor{fvdcyan}{HTML}{20BFC6}
\definecolor{fvdcyanlight}{HTML}{EFFCFB}
\definecolor{fvdmagenta}{HTML}{F10D69}
\definecolor{fvdmagentalight}{HTML}{FFF1F7}
\definecolor{fvdtext}{HTML}{163044}





%=================================================
% Pedagogische omgevingen
% PDF: tcolorbox
% HTML: learning-box
%=================================================

\newenvironment{voorkennis}
{%
    \ifdefined\HCode
        \HCode{%
<section class="learning-box learning-box--prior">
<div class="learning-box__header">
<span class="learning-box__icon" aria-hidden="true"></span>
<span class="learning-box__title">Voorkennis</span>
</div>
<div class="learning-box__content">
        }%
        \begin{itemize}
    \else
       \begin{tcolorbox}[
    enhanced,
    breakable,
    colback=fvdcyanlight,
    colframe=fvdcyan,
    coltext=fvdtext,
    boxrule=0.5pt,
    leftrule=4pt,
    arc=4mm,
    outer arc=4mm,
    left=7mm,
    right=6mm,
    top=5mm,
    bottom=5mm,
    before skip=6mm,
    after skip=6mm,
    title=Voorkennis,
    coltitle=fvdcyan,
    fonttitle=\bfseries\Large,
    attach title to upper={\par\smallskip},
    boxed title style={
        empty,
        left=0mm,
        right=0mm,
        top=0mm,
        bottom=0mm
    },
    overlay={
        \draw[
            fvdcyan,
            line width=0.5pt
        ]
        ([xshift=42mm,yshift=-13mm]frame.north west)
        --
        ([xshift=-7mm,yshift=-13mm]frame.north east);
    }
]
        \begin{itemize}
    \fi
}
{%
    \end{itemize}
    \ifdefined\HCode
        \HCode{%
</div>
</section>
        }%
    \else
        \end{tcolorbox}
    \fi
}


\newenvironment{leerdoelen}
{%
    \ifdefined\HCode
        \HCode{%
<section class="learning-box learning-box--goals">
<div class="learning-box__header">
<span class="learning-box__icon" aria-hidden="true"></span>
<span class="learning-box__title">Leerdoelen</span>
</div>
<div class="learning-box__content">
        }%
        \begin{itemize}
    \else
\begin{tcolorbox}[
    enhanced,
    breakable,
    colback=fvdmagentalight,
    colframe=fvdmagenta,
    coltext=fvdtext,
    boxrule=0.5pt,
    leftrule=4pt,
    arc=4mm,
    outer arc=4mm,
    left=7mm,
    right=6mm,
    top=5mm,
    bottom=5mm,
    before skip=6mm,
    after skip=6mm,
    title=Leerdoelen,
    coltitle=fvdmagenta,
    fonttitle=\bfseries\Large,
    attach title to upper={\par\smallskip},
    boxed title style={
        empty,
        left=0mm,
        right=0mm,
        top=0mm,
        bottom=0mm
    },
    overlay={
        \draw[
            fvdmagenta,
            line width=0.5pt
        ]
        ([xshift=42mm,yshift=-13mm]frame.north west)
        --
        ([xshift=-7mm,yshift=-13mm]frame.north east);
    }
]
        \begin{itemize}
    \fi
}
{%
    \end{itemize}
    \ifdefined\HCode
        \HCode{%
</div>
</section>
        }%
    \else
        \end{tcolorbox}
    \fi
}


\addPrintStyle{../../..}

\title{Elektrisch veld en potentiaal}
\author{Ingmar Herreman}

\begin{abstract}

  Een elektrische lading kan een andere lading aantrekken of afstoten,
  zelfs wanneer beide voorwerpen elkaar niet raken. In het vorige
  hoofdstuk beschreven we deze wisselwerking met de wet van Coulomb.

  Maar hoe kan een lading op afstand een kracht uitoefenen?

  Om die vraag te beantwoorden voeren we het begrip
  \textbf{elektrisch veld} in. Een geladen voorwerp verandert de
  elektrische eigenschappen van de ruimte rondom zich. Wanneer we een
  andere lading in die ruimte plaatsen, ondervindt ze een elektrische
  kracht.

  We onderzoeken hoe we zo'n elektrisch veld kunnen beschrijven met de
  elektrische veldvector en hoe veldlijnen ons helpen om de structuur
  van een veld zichtbaar te maken. We bestuderen radiale velden,
  homogene velden en het veld van een elektrische dipool. Wanneer
  meerdere ladingen aanwezig zijn, combineren we hun bijdragen met het
  superpositiebeginsel.

  Daarna bekijken we dezelfde elektrische wisselwerking vanuit een
  energetisch perspectief. We voeren elektrische potentiële energie,
  elektrische potentiaal en spanning in en onderzoeken het verband
  tussen veld, arbeid en energie.

  Zo ontstaan twee complementaire manieren om elektrische
  wisselwerkingen te beschrijven:

  \[
  \boxed{
  \text{kracht en veld}
  \qquad\longleftrightarrow\qquad
  \text{energie en potentiaal}
  }
  \]

  Ten slotte brengen we deze ideeën opnieuw samen met de mechanica.
  Een geladen deeltje in een elektrisch veld ondervindt een kracht en
  kan daardoor versnellen. Elektrische velden kunnen dus niet alleen
  krachten uitoefenen, maar ook de beweging en energie van geladen
  deeltjes veranderen.

\end{abstract}

\begin{document}

% FVD_CHAPTER_DATA_START
% Wordt automatisch aangepast vanuit index.tex
\fvdchapterdata{3}{Krachten op afstand}{2}
% FVD_CHAPTER_DATA_END

\maketitle


% ============================================================
% VOORKENNIS
% ============================================================

\begin{voorkennis}
  \item Je kent positieve en negatieve elektrische lading.
  \item Je kan de wet van Coulomb toepassen.
  \item Je kan elektrische krachten als vectoren voorstellen.
  \item Je kan meerdere vectoren optellen en ontbinden in componenten.
  \item Je kent de wetten van Newton.
  \item Je kent het gravitatieveld en de relatie
        $\vec F_g=m\vec g$.
  \item Je kent arbeid, energie en potentiële energie.
  \item Je kan werken met machten van tien en SI-eenheden.
\end{voorkennis}


% ============================================================
% LEERDOELEN
% ============================================================

\begin{leerdoelen}
  \item Je kan uitleggen waarom het veldmodel nuttig is om
        elektrische wisselwerkingen te beschrijven.
  \item Je kan de elektrische veldsterkte definiëren en interpreteren.
  \item Je kan het elektrisch veld van een puntlading bepalen.
  \item Je kan richting en zin van een elektrische veldvector bepalen.
  \item Je kan elektrische veldlijnen tekenen en interpreteren.
  \item Je kan radiale, homogene en dipoolvelden herkennen en
        analyseren.
  \item Je kan de veldsterkte vergelijken aan de hand van de dichtheid
        van veldlijnen.
  \item Je kan elektrische velden van meerdere bronladingen vectorieel
        combineren.
  \item Je kan de kracht op een positieve of negatieve lading in een
        elektrisch veld bepalen.
  \item Je kan elektrische potentiële energie interpreteren.
  \item Je kan elektrische potentiaal definiëren en berekenen.
  \item Je kan het onderscheid tussen potentiaal en spanning uitleggen.
  \item Je kan de potentiaal van meerdere puntladingen bepalen.
  \item Je kan het verband tussen elektrisch veld en potentiaal
        analyseren.
  \item Je kan de beweging en energie van geladen deeltjes in een
        homogeen elektrisch veld analyseren.
\end{leerdoelen}


% ============================================================
% 1. VAN KRACHT NAAR VELD
% ============================================================

\FVDsection{Van kracht naar veld}

In het vorige hoofdstuk onderzochten we de elektrische kracht tussen
twee puntladingen.

Voor een bronlading $Q$ en een tweede lading $q$ op afstand $r$ geldt:

\[
F_e
=
k\frac{|Qq|}{r^2}.
\]

De kracht ontstaat wanneer twee ladingen met elkaar interageren.

Maar we kunnen een andere vraag stellen.

Stel dat alleen de bronlading $Q$ aanwezig is.

Heeft het dan betekenis om iets te zeggen over de ruimte rondom die
lading?

Het antwoord is ja.

We zeggen dat een elektrische lading een
\textbf{elektrisch veld} rondom zich veroorzaakt.

Wanneer we vervolgens een tweede lading in dat veld plaatsen,
ondervindt die lading een elektrische kracht.

We kunnen de situatie dus voorstellen als:

\[
\boxed{
Q
\longrightarrow
\vec E
\longrightarrow
\vec F_e
}
\]

De bronlading $Q$ veroorzaakt het elektrische veld $\vec E$.

Een lading $q$ die in dat veld geplaatst wordt, ondervindt daardoor
een elektrische kracht $\vec F_e$.


% ============================================================
% 2. ANALOGIE MET GRAVITATIE
% ============================================================

\FVDsection{Een vertrouwde gedachte: het gravitatieveld}

Het veldbegrip is niet volledig nieuw.

Bij gravitatie schreven we:

\[
\boxed{
\vec F_g=m\vec g.
}
\]

De aarde veroorzaakt een gravitatieveld $\vec g$.

Een massa $m$ die zich in dat veld bevindt, ondervindt een
gravitatiekracht.

Schematisch:

\[
M
\longrightarrow
\vec g
\longrightarrow
\vec F_g.
\]

Bij elektrische wisselwerkingen bouwen we een volledig analoge
beschrijving op:

\[
Q
\longrightarrow
\vec E
\longrightarrow
\vec F_e.
\]

Dat geeft:

\[
\boxed{
\begin{array}{c|c}
\text{gravitatie} & \text{elektriciteit}\\
\hline
M & Q\\
m & q\\
\vec g & \vec E\\
\vec F_g=m\vec g & \vec F_e=q\vec E
\end{array}
}
\]


De analogie is krachtig, maar niet volledig.

Massa heeft in de klassieke mechanica één teken. Elektrische lading
kan positief of negatief zijn.

Dat zal belangrijke gevolgen hebben voor de richting van de kracht.


% ============================================================
% 3. TESTLADING
% ============================================================

\FVDsection{Hoe onderzoeken we een elektrisch veld?}

Om te onderzoeken wat het elektrische veld op een bepaalde plaats
doet, kunnen we ons voorstellen dat we daar een zeer kleine positieve
lading plaatsen.

Zo'n denkbeeldige lading noemen we een \textbf{testlading}.

We kiezen de testlading:

\begin{itemize}
  \item positief, zodat de richting van het veld ondubbelzinnig kan
        worden gedefinieerd;
  \item voldoende klein, zodat ze de oorspronkelijke ladingsverdeling
        zo weinig mogelijk verstoort.
\end{itemize}


Stel dat een positieve testlading $q_t$ op een bepaalde plaats een
kracht $\vec F_e$ ondervindt.

Dan gebruiken we die kracht om het elektrische veld op die plaats te
definiëren.


% ============================================================
% 4. DEFINITIE ELEKTRISCHE VELDSTERKTE
% ============================================================

\FVDsection{Elektrische veldsterkte}

\begin{definitie}[Elektrische veldsterkte]

De elektrische veldsterkte $\vec E$ in een punt is de elektrische
kracht per eenheid positieve testlading:

\[
\boxed{
\vec E
=
\frac{\vec F_e}{q_t}
}
\]

waarbij $q_t>0$ een positieve testlading is.

\end{definitie}


De elektrische veldsterkte is een \textbf{vectoriële grootheid}.

Ze heeft dus:

\begin{itemize}
  \item een grootte;
  \item een richting;
  \item een zin.
\end{itemize}


De richting en zin van $\vec E$ worden gedefinieerd als de richting en
zin van de kracht op een \textbf{positieve testlading}.


% ============================================================
% 5. EENHEID
% ============================================================

\FVDsection{Eenheid van elektrische veldsterkte}

Uit:

\[
E=\frac{F_e}{q}
\]

volgt:

\[
[E]
=
\frac{\mathrm N}{\mathrm C}.
\]

Dus:

\[
\boxed{
[E]=\mathrm{\frac NC}.
}
\]

Een elektrisch veld met grootte:

\[
E=500\,\mathrm{\frac NC}
\]

betekent dat een positieve testlading van $1\,\mathrm C$ op die plaats
een elektrische kracht van $500\,\mathrm N$ zou ondervinden.

Een testlading van $1\,\mathrm C$ is in veel praktische situaties
uitzonderlijk groot. De definitie blijft echter geldig voor kleinere
testladingen.


% ============================================================
% 6. VAN COULOMB NAAR VELD
% ============================================================

\FVDsection{Het elektrische veld van een puntlading}

We kunnen de veldsterkte van een puntlading afleiden uit de wet van
Coulomb.

Een bronlading $Q$ oefent op een positieve testlading $q_t$ een kracht
uit met grootte:

\[
F_e
=
k\frac{|Qq_t|}{r^2}.
\]

Volgens de definitie:

\[
E=\frac{F_e}{q_t}.
\]

Invullen geeft:

\[
E
=
\frac{1}{q_t}
k\frac{|Qq_t|}{r^2}.
\]

Omdat $q_t>0$:

\[
\boxed{
E
=
k\frac{|Q|}{r^2}.
}
\]


Dit resultaat is fundamenteel.

De testlading is verdwenen uit de formule.

De elektrische veldsterkte wordt bepaald door:

\begin{itemize}
  \item de bronlading $Q$;
  \item de afstand $r$ tot de bronlading.
\end{itemize}

Niet door de testlading waarmee we het veld onderzoeken.


% ============================================================
% 7. BRONLADING VERSUS TESTLADING
% ============================================================

\FVDsection{Bronlading en testlading niet verwarren}

Er spelen twee verschillende rollen.

De \textbf{bronlading} veroorzaakt het veld:

\[
Q
\longrightarrow
\vec E.
\]

De \textbf{testlading} ondervindt een kracht door dat veld:

\[
q
\longrightarrow
\vec F_e.
\]

Daarom:

\[
\boxed{
E=k\frac{|Q|}{r^2}
}
\]

maar:

\[
\boxed{
\vec F_e=q\vec E.
}
\]


\begin{opmerking}

Wanneer de testlading verdubbelt, verandert het elektrische veld niet.

De kracht op die testlading verdubbelt wel.

\end{opmerking}


\begin{oefening}[Veld of kracht?]{\oefB\ESS}

Een positieve testlading $q$ bevindt zich in een elektrisch veld met
veldsterkte $E$.

De testlading wordt vervangen door een lading $2q$ op dezelfde plaats.

Bepaal wat er gebeurt met:

\begin{enumerate}
  \item de elektrische veldsterkte;
  \item de elektrische kracht.
\end{enumerate}

Verklaar je antwoord.

\end{oefening}


% ============================================================
% 8. RICHTING VAN HET VELD
% ============================================================

\FVDsection{Richting en zin van het veld}

De richting van $\vec E$ is de richting waarin een positieve
testlading zou worden geduwd.


Voor een positieve bronlading:

\[
\boxed{
Q>0
\quad\Rightarrow\quad
\vec E\text{ wijst van }Q\text{ weg}.
}
\]


Voor een negatieve bronlading:

\[
\boxed{
Q<0
\quad\Rightarrow\quad
\vec E\text{ wijst naar }Q\text{ toe}.
}
\]


Dat volgt rechtstreeks uit aantrekking en afstoting.

Een positieve testlading wordt door een positieve bronlading afgestoten
en door een negatieve bronlading aangetrokken.


% ============================================================
% 9. RADIAAL VELD
% ============================================================

\FVDsection{Het radiale elektrische veld}

Het elektrische veld van een geïsoleerde puntlading is een
\textbf{radiaal veld}.

De veldvector ligt in elk punt op de rechte door het punt en de
bronlading.

Voor de grootte geldt:

\[
\boxed{
E(r)
=
k\frac{|Q|}{r^2}.
}
\]

Dus:

\[
\boxed{
E\propto\frac{1}{r^2}.
}
\]

Het elektrische veld van een puntlading volgt dezelfde
inverse-kwadratenwet als de Coulombkracht.


Wanneer:

\[
r\rightarrow2r,
\]

dan:

\[
E\rightarrow\frac14E.
\]

Wanneer:

\[
r\rightarrow3r,
\]

dan:

\[
E\rightarrow\frac19E.
\]


\begin{oefening}[Afstand en veldsterkte]{\oefB\ESS}

Op afstand $r$ van een puntlading bedraagt de elektrische veldsterkte
$E$.

Bepaal de veldsterkte op afstand:

\begin{enumerate}
  \item $2r$;
  \item $4r$;
  \item $r/2$.
\end{enumerate}

Werk met verhoudingen en gebruik de constante $k$ niet.

\end{oefening}


% ============================================================
% 10. VOORBEELD
% ============================================================

\begin{voorbeeld}[Veld van een puntlading]

Een puntlading:

\[
Q=+3{,}0\,\mathrm{\mu C}
\]

bevindt zich in de oorsprong.

Bereken de elektrische veldsterkte op een afstand:

\[
r=0{,}20\,\mathrm m.
\]

We gebruiken:

\[
E
=
k\frac{|Q|}{r^2}.
\]

Dus:

\[
E
=
8{,}99\times10^9
\frac{
3{,}0\times10^{-6}
}{
(0{,}20)^2
}.
\]

Daaruit volgt:

\[
\boxed{
E\approx6{,}7\times10^5\,\mathrm{\frac NC}.
}
\]

Omdat de bronlading positief is, wijst de veldvector van de bronlading
weg.

\end{voorbeeld}


% ============================================================
% 11. VELDLIJNEN
% ============================================================

\FVDsection{Een onzichtbaar veld zichtbaar maken}

Een elektrisch veld bestaat in elk punt van de ruimte.

Het is onmogelijk om in een tekening op elk punt een veldvector te
plaatsen.

Daarom gebruiken we \textbf{elektrische veldlijnen}.


\begin{definitie}[Veldlijn]

Een elektrische veldlijn is een denkbeeldige lijn waarvan de raaklijn
in elk punt de richting van de elektrische veldvector aangeeft.

\end{definitie}


Voor veldlijnen gelden enkele belangrijke regels:

\begin{itemize}
  \item veldlijnen vertrekken bij positieve ladingen;
  \item veldlijnen eindigen bij negatieve ladingen of lopen naar
        oneindig;
  \item de raaklijn aan een veldlijn geeft de richting van $\vec E$;
  \item de pijlen geven de zin van $\vec E$;
  \item veldlijnen kruisen elkaar nooit;
  \item waar de veldlijnen dichter bij elkaar liggen, is het veld
        sterker.
\end{itemize}


% ============================================================
% 12. WAAROM KRUISEN VELDLIJNEN NIET?
% ============================================================

\FVDsection{Waarom kunnen veldlijnen elkaar niet kruisen?}

Stel dat twee veldlijnen elkaar zouden kruisen.

In het snijpunt zouden dan twee verschillende raaklijnen bestaan.

Dat zou betekenen dat de elektrische veldvector in hetzelfde punt twee
verschillende richtingen heeft.

Maar een vector in één punt heeft één welbepaalde richting.

Daarom:

\[
\boxed{
\text{elektrische veldlijnen kunnen elkaar niet kruisen}.
}
\]


\begin{oefening}[Redeneren met veldlijnen]{\oefU\ESS}

Een leerling tekent twee elektrische veldlijnen die elkaar kruisen.

Leg zonder berekening uit waarom deze voorstelling fysisch onmogelijk
is.

\end{oefening}


% ============================================================
% 13. VELDLIJNDICHTHEID
% ============================================================

\FVDsection{Veldlijndichtheid en veldsterkte}

Een veldlijnendiagram geeft niet alleen informatie over de richting
van het veld.

Ook de dichtheid van de lijnen bevat informatie.

Waar de lijnen dichter bij elkaar liggen, stellen we een grotere
veldsterkte voor.

Waar ze verder uit elkaar liggen, stellen we een kleinere
veldsterkte voor.

Bij een puntlading spreiden de veldlijnen zich steeds verder uit
wanneer de afstand toeneemt.

Dat past bij:

\[
E\propto\frac{1}{r^2}.
\]


\begin{oefening}[Waar is het veld het sterkst?]{\oefB\ESS}

In een veldlijnendiagram liggen de veldlijnen in gebied A duidelijk
dichter bij elkaar dan in gebied B.

Vergelijk:

\[
E_A
\qquad\text{en}\qquad
E_B.
\]

Verklaar je antwoord.

\end{oefening}


% ============================================================
% 14. SUPERPOSITIE
% ============================================================

\FVDsection{Meerdere bronladingen}

Wanneer meerdere ladingen aanwezig zijn, veroorzaakt elke bronlading
haar eigen elektrisch veld.

In een bepaald punt geldt:

\[
\vec E_1,
\qquad
\vec E_2,
\qquad
\vec E_3,
\ldots
\]

De resulterende elektrische veldvector is de vectorsom:

\[
\boxed{
\vec E_\mathrm{res}
=
\vec E_1+\vec E_2+\vec E_3+\cdots
}
\]

Dit is opnieuw het \textbf{superpositiebeginsel}.


Let op het verschil met de Coulombkracht.

Bij het berekenen van het veld in een punt hebben we nog geen
testlading nodig.

We bepalen eerst:

\[
\vec E_\mathrm{res}.
\]

Pas wanneer een lading $q$ in dat punt wordt geplaatst, berekenen we:

\[
\vec F_e=q\vec E_\mathrm{res}.
\]


% ============================================================
% 15. SYSTEMATISCHE AANPAK
% ============================================================

\FVDsection{Een elektrisch veld systematisch bepalen}

Bij meerdere bronladingen gebruiken we de volgende strategie:

\[
\boxed{
\begin{array}{c}
\text{maak een schets}
\\
\downarrow
\\
\text{kies het punt waarin }\vec E\text{ gevraagd wordt}
\\
\downarrow
\\
\text{bepaal voor elke bronlading richting en zin}
\\
\downarrow
\\
\text{bereken elke veldsterkte afzonderlijk}
\\
\downarrow
\\
\text{ontbind indien nodig in componenten}
\\
\downarrow
\\
\text{tel de veldvectoren op}
\\
\downarrow
\\
\text{controleer richting, eenheid en grootteorde}
\end{array}
}
\]


Begin dus niet onmiddellijk te rekenen.

De vectortekening bepaalt hoe de afzonderlijke bijdragen gecombineerd
moeten worden.


% ============================================================
% 16. TWEE GELIJKE LADINGEN
% ============================================================

\FVDsection{Symmetrie gebruiken}

Beschouw twee gelijke positieve puntladingen:

\[
+Q
\qquad\qquad
+Q.
\]

We onderzoeken het elektrische veld precies in het midden tussen beide
ladingen.

De twee bronladingen bevinden zich op dezelfde afstand van het
middelpunt.

Dus:

\[
E_1=E_2.
\]

Maar hun veldvectoren hebben tegengestelde zin.

Daarom:

\[
\boxed{
\vec E_\mathrm{res}=\vec 0.
}
\]


Dit resultaat kunnen we vinden zonder één getal in te vullen.

Symmetrie kan dus een krachtiger hulpmiddel zijn dan rekenen.


\begin{oefening}[Symmetrie]{\oefU\ESS}

Twee identieke negatieve puntladingen bevinden zich op afstand $d$ van
elkaar.

Bepaal het elektrische veld precies in het midden tussen beide
ladingen.

Verklaar je antwoord met een vectortekening.

\end{oefening}


% ============================================================
% 17. ELEKTRISCHE DIPOOL
% ============================================================

\FVDsection{Het elektrische dipoolveld}

Een belangrijke configuratie bestaat uit twee even grote maar
tegengestelde ladingen:

\[
+Q
\qquad\qquad
-Q.
\]

Zo'n configuratie noemen we een \textbf{elektrische dipool}.

De veldlijnen vertrekken bij de positieve lading en eindigen bij de
negatieve lading.

Het veld is niet radiaal rond één enkel punt.

De richting en grootte van $\vec E$ veranderen van plaats tot plaats.

Het dipoolveld is belangrijk omdat dipoolachtige ladingsverdelingen
voorkomen in onder andere moleculen en gepolariseerde materie.


% ============================================================
% 18. INFLUENTIE OPNIEUW BEKEKEN
% ============================================================

\FVDsection{Influente werking in veldtaal}

In het vorige hoofdstuk onderzochten we elektrostatische influentie.

We kunnen dat verschijnsel nu preciezer beschrijven.

Een geladen voorwerp veroorzaakt een elektrisch veld.

Wanneer een geleider in dat veld wordt geplaatst, oefenen elektrische
krachten op de mobiele elektronen in de geleider.

Daardoor herverdelen de elektronen zich.

Influente werking kan dus worden beschreven als:

\[
\boxed{
\text{extern elektrisch veld}
\longrightarrow
\text{krachten op mobiele ladingen}
\longrightarrow
\text{ladingsherverdeling}.
}
\]

We hebben daarvoor geen rechtstreeks contact tussen beide voorwerpen
nodig.


% ============================================================
% 19. HOMOGEEN ELEKTRISCH VELD
% ============================================================

\FVDsection{Een homogeen elektrisch veld}

Niet elk elektrisch veld is radiaal.

Beschouw twee grote evenwijdige platen met tegengestelde lading.

Ver van de randen kunnen we het elektrische veld tussen de platen
benaderen als homogeen.

In een homogeen elektrisch veld is:

\[
\boxed{
\vec E=\text{constant}.
}
\]

Dat betekent dat zowel:

\begin{itemize}
  \item de grootte;
  \item de richting;
  \item de zin
\end{itemize}

van de veldvector overal dezelfde zijn.


De veldlijnen zijn dan:

\begin{itemize}
  \item recht;
  \item evenwijdig;
  \item gelijkmatig verdeeld.
\end{itemize}


Ze lopen van de positieve naar de negatieve plaat.


% ============================================================
% 20. RADIAAL VERSUS HOMOGEEN
% ============================================================

\FVDsection{Radiaal en homogeen veld vergelijken}

\[
\boxed{
\begin{array}{c|c}
\text{radiaal veld} & \text{homogeen veld}\\
\hline
\text{puntlading als bron}
&
\text{ideaal tussen evenwijdige platen}
\\[1mm]
E\propto1/r^2
&
E=\text{constant}
\\[1mm]
\text{richting verandert met plaats}
&
\text{richting overal gelijk}
\\[1mm]
\text{veldlijnen spreiden}
&
\text{veldlijnen evenwijdig}
\end{array}
}
\]


\begin{oefening}[Herken het veld]{\oefU\ESS}

Een veldlijnendiagram toont rechte, evenwijdige en gelijkmatig
verdeelde veldlijnen.

\begin{enumerate}
  \item Welk soort elektrisch veld wordt voorgesteld?
  \item Wat kan je zeggen over de grootte van $E$ op verschillende
        plaatsen?
  \item Wat kan je zeggen over de richting van $\vec E$?
\end{enumerate}

\end{oefening}


% ============================================================
% 21. KRACHT IN EEN ELEKTRISCH VELD
% ============================================================

\FVDsection{Van veld terug naar kracht}

Wanneer een lading $q$ in een elektrisch veld wordt geplaatst, geldt:

\[
\boxed{
\vec F_e=q\vec E.
}
\]

Dit lijkt sterk op:

\[
\vec F_g=m\vec g.
\]

Maar het teken van $q$ maakt een belangrijk verschil.


Voor:

\[
q>0
\]

hebben $\vec F_e$ en $\vec E$ dezelfde zin.


Voor:

\[
q<0
\]

hebben $\vec F_e$ en $\vec E$ tegengestelde zin.


\begin{opmerking}

Een negatieve lading verandert de richting van het elektrische veld
niet.

Het veld wordt bepaald door de bronladingen.

Het teken van de lading bepaalt wel de zin van de kracht die zij in dat
veld ondervindt.

\end{opmerking}


\begin{oefening}[Positief of negatief?]{\oefB\ESS}

Een elektrisch veld wijst horizontaal naar rechts.

Teken de elektrische kracht op:

\begin{enumerate}
  \item een positieve lading;
  \item een negatieve lading.
\end{enumerate}

\end{oefening}


% ============================================================
% 22. VOORBEELD KRACHT
% ============================================================

\begin{voorbeeld}[Elektron in een elektrisch veld]

Een elektron bevindt zich in een homogeen elektrisch veld:

\[
E=2{,}0\times10^4\,\mathrm{\frac NC}.
\]

De grootte van de kracht is:

\[
F_e=|q|E.
\]

Dus:

\[
F_e
=
(1{,}602\times10^{-19})
(2{,}0\times10^4).
\]

Daaruit volgt:

\[
\boxed{
F_e\approx3{,}2\times10^{-15}\,\mathrm N.
}
\]

Omdat het elektron negatief geladen is, wijst de kracht
\textbf{tegengesteld aan} de elektrische veldvector.

\end{voorbeeld}


% ============================================================
% 23. VAN KRACHT NAAR ENERGIE
% ============================================================

\FVDsection{Een tweede manier om elektrische interacties te beschrijven}

Tot nu toe beschreven we elektrische interacties met krachten en
velden.

Maar uit de mechanica kennen we een andere krachtige aanpak:

\[
\text{arbeid en energie}.
\]

Een elektrische kracht kan arbeid verrichten wanneer een geladen
deeltje zich verplaatst.

Daarom kunnen we aan de positie van een lading in een elektrisch veld
een \textbf{elektrische potentiële energie} koppelen.

Voor de arbeid van de elektrische kracht geldt:

\[
\boxed{
W_e=-\Delta E_{\mathrm{pot,e}}.
}
\]

Wanneer de elektrische kracht positieve arbeid verricht:

\[
\Delta E_{\mathrm{pot,e}}<0.
\]

De elektrische potentiële energie neemt dan af.


% ============================================================
% 24. POTENTIELE ENERGIE VAN TWEE PUNTLADINGEN
% ============================================================

\FVDsection{Elektrische potentiële energie van twee puntladingen}

Voor twee puntladingen $Q$ en $q$ op afstand $r$ definiëren we, met
het nulpunt op oneindig:

\[
\boxed{
E_{\mathrm{pot,e}}
=
k\frac{Qq}{r}.
}
\]

Let op: hier gebruiken we geen absolute waarde.

Het teken bevat fysische informatie.


Voor gelijknamige ladingen:

\[
Qq>0
\]

en dus:

\[
E_{\mathrm{pot,e}}>0.
\]


Voor tegengestelde ladingen:

\[
Qq<0
\]

en dus:

\[
E_{\mathrm{pot,e}}<0.
\]


% ============================================================
% 25. WAAROM EEN TEKEN?
% ============================================================

\FVDsection{Wat betekent het teken van potentiële energie?}

Het nulpunt van de potentiële energie kiezen we bij:

\[
r\rightarrow\infty.
\]

Voor twee gelijknamige ladingen moeten we arbeid leveren om ze dichter
bij elkaar te brengen tegen de afstotende elektrische kracht in.

De potentiële energie wordt positief.

Voor tegengestelde ladingen trekken de ladingen elkaar spontaan aan.

Wanneer ze dichter bij elkaar komen, kan de elektrische kracht
positieve arbeid verrichten.

De potentiële energie wordt dan negatiever.


\begin{oefening}[Teken van de energie]{\oefU\ESS}

Bepaal zonder numerieke berekening het teken van de elektrische
potentiële energie voor:

\begin{enumerate}
  \item twee positieve puntladingen;
  \item twee negatieve puntladingen;
  \item een positieve en een negatieve puntlading.
\end{enumerate}

Verklaar telkens fysisch.

\end{oefening}


% ============================================================
% 26. VAN POTENTIELE ENERGIE NAAR POTENTIAAL
% ============================================================

\FVDsection{Van energie naar potentiaal}

De elektrische potentiële energie:

\[
E_{\mathrm{pot,e}}
=
k\frac{Qq}{r}
\]

hangt af van zowel de bronlading $Q$ als de lading $q$ die we in het
veld plaatsen.

Maar net zoals we bij kracht het veld wilden beschrijven onafhankelijk
van de testlading, willen we dat ook bij energie.

Daarom definiëren we de \textbf{elektrische potentiaal}.


\begin{definitie}[Elektrische potentiaal]

De elektrische potentiaal $V$ in een punt is de elektrische potentiële
energie per eenheid lading:

\[
\boxed{
V
=
\frac{E_{\mathrm{pot,e}}}{q}.
}
\]

\end{definitie}


Voor een puntlading volgt:

\[
V
=
\frac{1}{q}
k\frac{Qq}{r}.
\]

Dus:

\[
\boxed{
V
=
k\frac{Q}{r}.
}
\]


% ============================================================
% 27. VELD EN POTENTIAAL
% ============================================================

\FVDsection{Veld en potentiaal: twee beschrijvingen}

Voor een puntlading:

\[
\boxed{
E
=
k\frac{|Q|}{r^2}
}
\]

en:

\[
\boxed{
V
=
k\frac{Q}{r}.
}
\]

Er zijn belangrijke verschillen.

De elektrische veldsterkte $\vec E$ is een vector.

De elektrische potentiaal $V$ is een scalair.

Bovendien:

\[
E\propto\frac1{r^2},
\]

terwijl:

\[
V\propto\frac1r.
\]


\begin{center}
\[
\boxed{
\begin{array}{c|c}
\text{elektrisch veld} & \text{elektrische potentiaal}\\
\hline
\vec E & V\\
\text{vector} & \text{scalair}\\
\mathrm{N/C} & \mathrm V\\
E\propto1/r^2 & V\propto1/r
\end{array}
}
\]
\end{center}


% ============================================================
% 28. EENHEID VAN POTENTIAAL
% ============================================================

\FVDsection{De volt}

Uit:

\[
V=\frac{E_{\mathrm{pot,e}}}{q}
\]

volgt:

\[
[V]=\frac{\mathrm J}{\mathrm C}.
\]

Deze eenheid noemen we de \textbf{volt}:

\[
\boxed{
1\,\mathrm V
=
1\,\frac{\mathrm J}{\mathrm C}.
}
\]

Een potentiaalverschil van één volt betekent dus een verschil van één
joule potentiële energie per coulomb lading.


% ============================================================
% 29. POTENTIAAL IS GEEN POTENTIELE ENERGIE
% ============================================================

\FVDsection{Potentiaal is geen energie}

De begrippen worden gemakkelijk verward.

Potentiaal $V$ is een eigenschap van een plaats in een elektrisch veld.

Potentiële energie $E_{\mathrm{pot,e}}$ hangt daarnaast af van de
lading die op die plaats aanwezig is.

Hun verband is:

\[
\boxed{
E_{\mathrm{pot,e}}=qV.
}
\]


Op dezelfde plaats hebben verschillende ladingen dus dezelfde
potentiaal $V$, maar niet noodzakelijk dezelfde potentiële energie.


\begin{oefening}[Zelfde plaats, andere lading]{\oefB\ESS}

Op een bepaalde plaats is:

\[
V=250\,\mathrm V.
\]

Bereken de elektrische potentiële energie van:

\begin{enumerate}
  \item $q=+2{,}0\,\mathrm{\mu C}$;
  \item $q=-2{,}0\,\mathrm{\mu C}$.
\end{enumerate}

Vergelijk beide resultaten.

\end{oefening}


% ============================================================
% 30. SUPERPOSITIE VAN POTENTIAAL
% ============================================================

\FVDsection{Potentiaal van meerdere bronladingen}

Wanneer meerdere bronladingen aanwezig zijn, draagt elke bronlading bij
aan de potentiaal.

Voor één puntlading:

\[
V_i
=
k\frac{Q_i}{r_i}.
\]

Omdat potentiaal een scalaire grootheid is, geldt:

\[
\boxed{
V_\mathrm{tot}
=
V_1+V_2+V_3+\cdots
}
\]

of:

\[
\boxed{
V
=
k\sum_i\frac{Q_i}{r_i}.
}
\]


Hier moeten de tekens van de ladingen wel behouden blijven.

Een positieve bronlading levert een positieve bijdrage.

Een negatieve bronlading levert een negatieve bijdrage.


% ============================================================
% 31. VECTOR VERSUS SCALAIR
% ============================================================

\FVDsection{Waarom is potentiaal soms eenvoudiger?}

Voor het elektrische veld:

\[
\boxed{
\vec E_\mathrm{res}
=
\sum_i\vec E_i.
}
\]

Daarvoor moeten we rekening houden met richting en zin.

Voor potentiaal:

\[
\boxed{
V_\mathrm{tot}
=
\sum_iV_i.
}
\]

Dat is een algebraïsche som.

Dit verschil tussen vectoren en scalairen is een van de redenen waarom
potentiaal bij complexe elektrische configuraties zo'n nuttig begrip
is.


\begin{oefening}[Midden tussen twee ladingen]{\oefU\ESS}

Twee puntladingen:

\[
+Q
\qquad\text{en}\qquad
-Q
\]

bevinden zich op gelijke afstand van punt $P$.

\begin{enumerate}
  \item Bepaal de totale potentiaal in $P$.
  \item Mag je daaruit besluiten dat ook
        $\vec E(P)=\vec 0$?
  \item Onderzoek de richting van beide veldvectoren en formuleer een
        besluit.
\end{enumerate}

\end{oefening}


% ============================================================
% 32. SPANNING
% ============================================================

\FVDsection{Potentiaalverschil en spanning}

In toepassingen is vaak niet de potentiaal in één punt belangrijk,
maar het verschil tussen twee punten.

We definiëren het potentiaalverschil van A naar B als:

\[
\Delta V
=
V_B-V_A.
\]

Een potentiaalverschil noemen we ook een
\textbf{elektrische spanning}.

We kunnen dus schrijven:

\[
\boxed{
U_{AB}=V_A-V_B
}
\]

wanneer $U_{AB}$ de spanning van A ten opzichte van B voorstelt.

Let steeds op de gebruikte volgorde van de punten.


Voor de verandering van potentiële energie geldt:

\[
\boxed{
\Delta E_{\mathrm{pot,e}}
=
q\Delta V.
}
\]

Dus:

\[
\boxed{
E_{\mathrm{pot,e},B}
-
E_{\mathrm{pot,e},A}
=
q(V_B-V_A).
}
\]


% ============================================================
% 33. ARBEID EN SPANNING
% ============================================================

\FVDsection{Arbeid, energie en spanning}

Omdat:

\[
W_e
=
-\Delta E_{\mathrm{pot,e}},
\]

en:

\[
\Delta E_{\mathrm{pot,e}}
=
q(V_B-V_A),
\]

volgt:

\[
W_{e,A\rightarrow B}
=
q(V_A-V_B).
\]

Dus:

\[
\boxed{
W_{e,A\rightarrow B}
=
qU_{AB}.
}
\]

Een elektrisch potentiaalverschil kan dus rechtstreeks gekoppeld
worden aan de arbeid die de elektrische kracht kan verrichten.


% ============================================================
% 34. POSITIEVE EN NEGATIEVE LADING
% ============================================================

\FVDsection{Welke kant beweegt een lading spontaan uit?}

Een positieve lading die vrij kan bewegen, wordt versneld in de
richting van $\vec E$.

Ze beweegt daardoor spontaan naar lagere elektrische potentiaal.

Voor een positieve lading:

\[
q>0
\quad\Rightarrow\quad
E_{\mathrm{pot,e}}=qV.
\]

Lagere $V$ betekent dus lagere potentiële energie.


Voor een negatieve lading is de situatie omgekeerd.

Omdat:

\[
q<0,
\]

kan een stijging van $V$ juist overeenkomen met een daling van de
potentiële energie.

Daarom:

\[
\boxed{
\text{een negatieve lading wordt versneld tegengesteld aan }\vec E.
}
\]


\begin{oefening}[Elektron tussen twee punten]{\oefU\ESS}

Voor twee punten geldt:

\[
V_A=100\,\mathrm V,
\qquad
V_B=500\,\mathrm V.
\]

Een elektron kan vrij bewegen.

\begin{enumerate}
  \item In welke richting wijst het elektrische veld langs de
        verbindingslijn?
  \item Naar welk punt wordt het elektron versneld?
  \item Neemt zijn elektrische potentiële energie daarbij toe of af?
\end{enumerate}

Verklaar de drie antwoorden samenhangend.

\end{oefening}


% ============================================================
% 35. EQUIPOTENTIAAL
% ============================================================

\FVDsection{Equipotentiaallijnen en equipotentiaalvlakken}

Een \textbf{equipotentiaallijn} verbindt punten met dezelfde
elektrische potentiaal.

In drie dimensies spreken we van een
\textbf{equipotentiaalvlak}.

Langs zo'n lijn of vlak geldt:

\[
\Delta V=0.
\]

Dus:

\[
\Delta E_{\mathrm{pot,e}}=q\Delta V=0.
\]

De elektrische kracht verricht bijgevolg geen arbeid wanneer een
lading langs een equipotentiaal wordt verplaatst.


Elektrische veldlijnen staan loodrecht op equipotentiaallijnen:

\[
\boxed{
\vec E
\perp
\text{equipotentiaallijnen}.
}
\]


% ============================================================
% 36. WAAROM LOODRECHT?
% ============================================================

\FVDsection{Waarom staan veldlijnen loodrecht op equipotentialen?}

Stel dat de elektrische veldvector een component langs een
equipotentiaallijn zou hebben.

Dan zou de elektrische kracht op een positieve testlading eveneens een
component langs die lijn hebben.

Bij een kleine verplaatsing langs de equipotentiaal zou de elektrische
kracht dan arbeid verrichten.

Maar langs een equipotentiaal geldt:

\[
\Delta V=0
\]

en dus:

\[
W_e=0.
\]

Daarom kan de elektrische veldvector geen component langs de
equipotentiaal hebben.

Dus:

\[
\boxed{
\vec E\perp\text{equipotentiaal}.
}
\]


% ============================================================
% 37. HOMOGEEN VELD EN POTENTIAAL
% ============================================================

\FVDsection{Potentiaalverschil in een homogeen veld}

Beschouw een homogeen elektrisch veld met veldsterkte $E$.

We verplaatsen ons over een afstand $d$ in de richting van het veld.

De potentiaal neemt af.

Voor de grootte van het potentiaalverschil geldt:

\[
\boxed{
|\Delta V|=Ed.
}
\]

Daaruit:

\[
\boxed{
E=\frac{|\Delta V|}{d}.
}
\]


Dit geeft een tweede eenheid voor elektrische veldsterkte:

\[
[E]
=
\frac{\mathrm V}{\mathrm m}.
\]

Dus:

\[
\boxed{
1\,\mathrm{\frac NC}
=
1\,\mathrm{\frac Vm}.
}
\]


% ============================================================
% 38. VOORBEELD HOMOGEEN VELD
% ============================================================

\begin{voorbeeld}[Twee evenwijdige platen]

Twee evenwijdige platen bevinden zich op:

\[
d=5{,}0\,\mathrm{cm}
\]

van elkaar.

Het potentiaalverschil tussen de platen bedraagt:

\[
|\Delta V|=600\,\mathrm V.
\]

In het ideale model is:

\[
E=\frac{|\Delta V|}{d}.
\]

Met:

\[
d=0{,}050\,\mathrm m
\]

vinden we:

\[
E
=
\frac{600}{0{,}050}.
\]

Dus:

\[
\boxed{
E=1{,}2\times10^4\,\mathrm{\frac Vm}.
}
\]

Omdat:

\[
1\,\mathrm{\frac Vm}
=
1\,\mathrm{\frac NC},
\]

kunnen we ook schrijven:

\[
\boxed{
E=1{,}2\times10^4\,\mathrm{\frac NC}.
}
\]

\end{voorbeeld}


% ============================================================
% 39. VAN ELEKTRICITEIT TERUG NAAR MECHANICA
% ============================================================

\FVDsection{Een geladen deeltje in een elektrisch veld}

Nu kunnen we de elektrische leerstof verbinden met de wetten van
Newton.

Een geladen deeltje in een elektrisch veld ondervindt:

\[
\vec F_e=q\vec E.
\]

Volgens Newton II:

\[
\vec F_\mathrm{res}=m\vec a.
\]

Wanneer de elektrische kracht de enige relevante kracht is:

\[
m\vec a=q\vec E.
\]

Dus:

\[
\boxed{
\vec a
=
\frac{q}{m}\vec E.
}
\]


In een homogeen elektrisch veld zijn $q$, $m$ en $\vec E$ constant.

Daarom is ook:

\[
\boxed{
\vec a=\text{constant}.
}
\]

Een geladen deeltje voert dan een eenparig versnelde beweging uit in
de richting van de elektrische kracht.


% ============================================================
% 40. POSITIEF EN NEGATIEF DEELTJE
% ============================================================

\FVDsection{Welke kant versnelt het deeltje uit?}

Voor een positief geladen deeltje:

\[
q>0
\]

geldt:

\[
\vec a\parallel\vec E
\]

met dezelfde zin.


Voor een negatief geladen deeltje:

\[
q<0
\]

geldt:

\[
\vec a\parallel\vec E
\]

met tegengestelde zin.


De grootte van de versnelling is:

\[
\boxed{
a
=
\frac{|q|E}{m}.
}
\]


Daaruit zien we dat niet alleen de lading belangrijk is.

Ook de massa speelt een rol.

De verhouding:

\[
\boxed{
\frac{|q|}{m}
}
\]

bepaalt hoe sterk een deeltje op een gegeven elektrisch veld reageert.


% ============================================================
% 41. ENERGIE VAN EEN VERSNELD DEELTJE
% ============================================================

\FVDsection{Elektrische energie wordt kinetische energie}

Wanneer alleen de elektrische kracht arbeid verricht, geldt behoud van
mechanische energie:

\[
E_k+E_{\mathrm{pot,e}}=\text{constant}.
\]

Dus:

\[
\Delta E_k
=
-\Delta E_{\mathrm{pot,e}}.
\]

Omdat:

\[
\Delta E_{\mathrm{pot,e}}
=
q\Delta V,
\]

volgt:

\[
\boxed{
\Delta E_k
=
-q\Delta V.
}
\]

Of voor beweging van A naar B:

\[
\boxed{
\Delta E_k
=
q(V_A-V_B).
}
\]


Wanneer een deeltje vanuit rust versnelt:

\[
E_{k,A}=0,
\]

dan kan:

\[
\frac12mv^2
=
q(V_A-V_B)
\]

worden gebruikt, op voorwaarde dat het rechterlid positief is.


% ============================================================
% 42. VOORBEELD PROTON
% ============================================================

\begin{voorbeeld}[Proton versnellen]

Een proton vertrekt uit rust en beweegt door een potentiaalverschil
waarbij zijn elektrische potentiële energie afneemt met:

\[
1{,}60\times10^{-16}\,\mathrm J.
\]

De kinetische energie neemt met dezelfde hoeveelheid toe:

\[
\Delta E_k
=
1{,}60\times10^{-16}\,\mathrm J.
\]

Met:

\[
E_k=\frac12m_pv^2
\]

vinden we:

\[
v
=
\sqrt{
\frac{2E_k}{m_p}
}.
\]

Dit voorbeeld toont hoe een elektrisch potentiaalverschil rechtstreeks
kan worden gebruikt om geladen deeltjes te versnellen.

\end{voorbeeld}


% ============================================================
% 43. ELEKTRONVOLT
% ============================================================

\FVDsection{Een handige energie-eenheid: de elektronvolt}

Op atomaire en subatomaire schaal is de joule vaak een onhandig grote
eenheid.

Daarom gebruiken fysici vaak de \textbf{elektronvolt}.

Eén elektronvolt is de energieverandering van een elementaire lading
die een potentiaalverschil van één volt doorloopt:

\[
\boxed{
1\,\mathrm{eV}
=
1{,}602\times10^{-19}\,\mathrm J.
}
\]

Let op:

\[
\mathrm{eV}
\]

is een eenheid van \textbf{energie}, niet van spanning.


% ============================================================
% 44. AFBUIGING VAN EEN DEELTJESBUNDEL
% ============================================================

\FVDsection{STEM: geladen deeltjes afbuigen}

Beschouw een positief geladen deeltje dat horizontaal een gebied met
een homogeen verticaal elektrisch veld binnenkomt.

Horizontaal werkt geen elektrische kracht.

Daarom is:

\[
a_x=0.
\]

De horizontale beweging is eenparig:

\[
x=v_{0x}t.
\]


Verticaal geldt:

\[
F_y=qE.
\]

Dus:

\[
a_y=\frac{qE}{m}.
\]

Wanneer:

\[
v_{0y}=0,
\]

dan:

\[
y
=
\frac12a_yt^2.
\]

Dus:

\[
y
=
\frac12
\frac{qE}{m}
t^2.
\]

Uit:

\[
t=\frac{x}{v_{0x}}
\]

volgt:

\[
y
=
\frac12
\frac{qE}{m}
\frac{x^2}{v_{0x}^2}.
\]

Daarom:

\[
\boxed{
y
=
\frac{qE}{2mv_{0x}^2}x^2.
}
\]


De baan is dus een \textbf{parabool}.


% ============================================================
% 45. EEN VERTROUWDE PARABOOL
% ============================================================

\FVDsection{De elektrische versie van een worpbeweging}

De vorige beweging lijkt sterk op een horizontale worp.

Bij een horizontale worp:

\[
a_x=0,
\qquad
a_y=g.
\]

Bij een geladen deeltje in een homogeen elektrisch veld:

\[
a_x=0,
\qquad
a_y=\frac{qE}{m}.
\]

In beide gevallen combineren we:

\begin{itemize}
  \item een eenparige beweging in één richting;
  \item een eenparig versnelde beweging in de loodrechte richting.
\end{itemize}

Daardoor ontstaat in beide gevallen een parabolische baan.

Hetzelfde wiskundige model kan dus verschillende fysische situaties
beschrijven.


% ============================================================
% 46. MODEL EN WERKELIJKHEID
% ============================================================

\FVDsection{Hoe goed is ons model?}

Wanneer we schrijven:

\[
E=\frac{|\Delta V|}{d},
\]

voor twee evenwijdige platen, gebruiken we een ideaal model.

We veronderstellen onder andere:

\begin{itemize}
  \item grote evenwijdige platen;
  \item een afstand die klein is ten opzichte van hun afmetingen;
  \item een gebied voldoende ver van de randen;
  \item verwaarloosbare invloed van andere ladingen en velden.
\end{itemize}

Dicht bij de randen zijn de veldlijnen niet perfect recht en
evenwijdig.

Daar treedt een \textbf{randeffect} op.

Een fysisch model hoeft dus niet overal exact te zijn om zeer bruikbaar
te zijn.


% ============================================================
% 47. OEFENINGENBUNDEL
% ============================================================

\FVDsection{Oefeningen}


% ------------------------------------------------------------
% BASIS
% ------------------------------------------------------------

\FVDsubsection{Basis}


\begin{oefening}[Veldsterkte]{\oefB\ESS}

Een puntlading:

\[
Q=+5{,}0\,\mathrm{\mu C}
\]

bevindt zich in de oorsprong.

Bereken de elektrische veldsterkte op:

\[
r=0{,}30\,\mathrm m.
\]

Geef ook richting en zin.

\end{oefening}


\begin{oefening}[Negatieve bronlading]{\oefB\ESS}

Een puntlading:

\[
Q=-8{,}0\,\mathrm{nC}
\]

bevindt zich op afstand:

\[
r=4{,}0\,\mathrm{cm}
\]

van punt $P$.

Bereken $E$ in $P$ en beschrijf de zin van $\vec E$.

\end{oefening}


\begin{oefening}[Kracht uit veld]{\oefB\ESS}

Een positieve lading:

\[
q=2{,}0\,\mathrm{\mu C}
\]

bevindt zich in een homogeen elektrisch veld:

\[
E=3{,}0\times10^4\,\mathrm{\frac NC}.
\]

Bereken de elektrische kracht.

\end{oefening}


\begin{oefening}[Elektron]{\oefB\ESS}

Een elektron bevindt zich in hetzelfde veld als in de vorige
oefening.

\begin{enumerate}
  \item Bereken de grootte van de kracht.
  \item Vergelijk de zin van $\vec F_e$ met die van $\vec E$.
\end{enumerate}

\end{oefening}


\begin{oefening}[Potentiaal van een puntlading]{\oefB\ESS}

Een puntlading:

\[
Q=+4{,}0\,\mathrm{\mu C}
\]

bevindt zich op afstand:

\[
r=0{,}50\,\mathrm m
\]

van punt $P$.

Bereken de elektrische potentiaal in $P$ wanneer het nulpunt op
oneindig ligt.

\end{oefening}


\begin{oefening}[Potentiële energie]{\oefB\ESS}

Op een plaats met:

\[
V=-250\,\mathrm V
\]

wordt een lading geplaatst:

\[
q=+3{,}0\,\mathrm{\mu C}.
\]

Bereken de elektrische potentiële energie.

\end{oefening}


\begin{oefening}[Homogeen veld]{\oefB\ESS}

Tussen twee evenwijdige platen bedraagt het potentiaalverschil:

\[
900\,\mathrm V.
\]

De afstand tussen de platen is:

\[
3{,}0\,\mathrm{cm}.
\]

Bereken de elektrische veldsterkte in het ideale homogene gebied.

\end{oefening}


% ------------------------------------------------------------
% VERDIEPING
% ------------------------------------------------------------

\FVDsubsection{Verdieping}


\begin{oefening}[Veld en kracht onderscheiden]{\oefU\ESS}

Een leerling zegt:

``Als ik de testlading driemaal groter maak, wordt het elektrische veld
ook driemaal groter.''

Analyseer de uitspraak.

Gebruik zowel:

\[
\vec E=\frac{\vec F_e}{q}
\]

als de formule voor het veld van een puntlading in je verklaring.

\end{oefening}


\begin{oefening}[Twee positieve bronladingen]{\oefU\ESS}

Twee puntladingen:

\[
Q_1=+2{,}0\,\mathrm{\mu C},
\qquad
Q_2=+8{,}0\,\mathrm{\mu C}
\]

bevinden zich op een afstand:

\[
d=0{,}60\,\mathrm m
\]

van elkaar.

Bepaal op de verbindingslijn tussen beide ladingen het punt waar:

\[
\vec E_\mathrm{res}=\vec 0.
\]

Maak eerst een kwalitatieve voorspelling van de ligging.

\end{oefening}


\begin{oefening}[Veld nul, potentiaal niet nul]{\oefU\ESS}

Twee gelijke positieve ladingen bevinden zich symmetrisch rond punt
$P$.

\begin{enumerate}
  \item Bepaal kwalitatief $\vec E(P)$.
  \item Bepaal het teken van $V(P)$.
  \item Leg uit waarom een nulveld niet noodzakelijk een nulpotentiaal
        betekent.
\end{enumerate}

\end{oefening}


\begin{oefening}[Potentiaal nul, veld niet nul]{\oefU\ESS}

Een positieve en een even grote negatieve puntlading bevinden zich
symmetrisch rond punt $P$.

\begin{enumerate}
  \item Bepaal $V(P)$.
  \item Onderzoek de veldvectoren in $P$.
  \item Is $\vec E(P)$ nul?
  \item Leg uit waarom $V=0$ niet noodzakelijk betekent dat
        $\vec E=\vec 0$.
\end{enumerate}

\end{oefening}


\begin{oefening}[Drie bronladingen]{\oefU\ESS}

Drie puntladingen liggen op de $x$-as:

\[
Q_1=+2{,}0\,\mathrm{\mu C}
\quad\text{bij}\quad
x=0,
\]

\[
Q_2=-4{,}0\,\mathrm{\mu C}
\quad\text{bij}\quad
x=0{,}20\,\mathrm m,
\]

\[
Q_3=+3{,}0\,\mathrm{\mu C}
\quad\text{bij}\quad
x=0{,}50\,\mathrm m.
\]

Bepaal:

\begin{enumerate}
  \item de elektrische veldsterkte in
        $x=0{,}35\,\mathrm m$;
  \item de elektrische potentiaal in datzelfde punt.
\end{enumerate}

Vergelijk de werkwijze voor de vectoriële en scalaire grootheid.

\end{oefening}


\begin{oefening}[Energie en snelheid]{\oefU\ESS}

Een proton vertrekt uit rust en wordt versneld door een
potentiaalverschil van:

\[
500\,\mathrm V.
\]

Verwaarloos andere krachten.

Bereken de snelheid van het proton nadat het dit potentiaalverschil
heeft doorlopen.

Gebruik:

\[
m_p=1{,}67\times10^{-27}\,\mathrm{kg}.
\]

\end{oefening}


\begin{oefening}[Redeneren zonder rekenen]{\oefU\ESS}

Een positieve lading wordt losgelaten in een niet-homogeen elektrisch
veld.

Ze beweegt aanvankelijk in de richting van een gebied met lagere
potentiaal.

Bespreek wat je kan besluiten over:

\begin{enumerate}
  \item de richting van de elektrische kracht;
  \item de verandering van de potentiële energie;
  \item de verandering van de kinetische energie wanneer alleen de
        elektrische kracht arbeid verricht.
\end{enumerate}

\end{oefening}


% ------------------------------------------------------------
% GEVORDERD
% ------------------------------------------------------------

\FVDsubsection{Gevorderd}


\begin{oefening}[Veld in twee dimensies]{\oefG}

Twee positieve puntladingen:

\[
Q_1=2Q
\qquad\text{en}\qquad
Q_2=Q
\]

bevinden zich op de $x$-as.

Een punt $P$ bevindt zich boven het midden tussen beide ladingen.

Stel zelf geschikte coördinaten in en bepaal symbolisch:

\begin{enumerate}
  \item $\vec E_1(P)$;
  \item $\vec E_2(P)$;
  \item de componenten van $\vec E_\mathrm{res}(P)$;
  \item de richting van het resulterende veld.
\end{enumerate}

Controleer of je antwoord overeenkomt met wat je op basis van de
symmetrie en de ongelijke ladingen verwacht.

\end{oefening}


\begin{oefening}[Veld en potentiaal vergelijken]{\oefG}

Voor een puntlading geldt:

\[
E(r)=k\frac{|Q|}{r^2}
\]

en:

\[
V(r)=k\frac{Q}{r}.
\]

Onderzoek wiskundig en fysisch:

\begin{enumerate}
  \item hoe beide functies veranderen met $r$;
  \item welke functie het snelst naar nul gaat;
  \item waarom $E$ en $V$ verschillende machten van $r$ bevatten;
  \item welk verband je ziet tussen de helling van $V(r)$ en de
        grootte van $E(r)$.
\end{enumerate}

\end{oefening}


\begin{oefening}[Van potentiaal naar veld]{\oefG}

Langs de $x$-as wordt de elektrische potentiaal beschreven door:

\[
V(x)=\frac{A}{x}
\]

voor $x>0$, met $A>0$.

Voor leerlingen die vertrouwd zijn met afgeleiden geldt in één
dimensie:

\[
E_x=-\frac{dV}{dx}.
\]

\begin{enumerate}
  \item Bepaal $E_x(x)$.
  \item Vergelijk het resultaat met het radiale veld van een positieve
        puntlading.
  \item Verklaar de betekenis van het minteken.
\end{enumerate}

\end{oefening}


\begin{oefening}[Elektrische afbuiging]{\oefG}

Een positief geladen deeltje met massa $m$ en lading $q$ komt
horizontaal met snelheid $v_0$ een gebied binnen met een homogeen
verticaal elektrisch veld $E$.

Het veldgebied heeft horizontale lengte $L$.

\begin{enumerate}
  \item Bepaal hoe lang het deeltje zich tussen de platen bevindt.
  \item Bepaal de verticale versnelling.
  \item Bepaal de verticale uitwijking bij het verlaten van het
        veldgebied.
  \item Toon aan dat de baan tussen de platen parabolisch is.
  \item Onderzoek hoe de uitwijking verandert wanneer $v_0$ wordt
        verdubbeld.
\end{enumerate}

\end{oefening}


\begin{oefening}[Modelkritiek bij evenwijdige platen]{\oefG}

Een leerling gebruikt overal tussen en rond twee eindige geladen platen:

\[
E=\frac{|\Delta V|}{d}.
\]

Analyseer deze keuze.

Bespreek:

\begin{enumerate}
  \item onder welke omstandigheden de formule een goed model vormt;
  \item wat er bij de randen van de platen gebeurt;
  \item waarom een fysisch model bruikbaar kan zijn zonder exact de
        volledige werkelijkheid te beschrijven.
\end{enumerate}

\end{oefening}


% ============================================================
% 48. ICT-ONDERZOEK
% ============================================================

\FVDsection{Onderzoeken met ICT}

Gebruik GeoGebra, Desmos, een spreadsheet of een geschikte
veldsimulatie om het elektrische veld en de elektrische potentiaal van
puntladingen te onderzoeken.

Onderzoek eerst één positieve puntlading.

Maak grafieken van:

\[
E(r)=k\frac{|Q|}{r^2}
\]

en:

\[
V(r)=k\frac{Q}{r}.
\]

Vergelijk:

\begin{enumerate}
  \item de vorm van beide grafieken;
  \item hun gedrag voor grote $r$;
  \item hun gedrag wanneer $r$ kleiner wordt;
  \item de invloed van een verdubbeling van $Q$.
\end{enumerate}

Onderzoek daarna twee ladingen.

Vergelijk bijvoorbeeld:

\[
(+Q,+Q)
\]

met:

\[
(+Q,-Q).
\]

Maak voor verschillende punten een voorstelling van:

\[
\vec E
\]

en:

\[
V.
\]

Zoek punten waar:

\[
\vec E=\vec 0
\]

en punten waar:

\[
V=0.
\]

Onderzoek of dit dezelfde punten zijn.

Formuleer op basis van je waarnemingen een besluit.


% ============================================================
% 49. MISCONCEPTIES
% ============================================================

\FVDsection{Veelvoorkomende denkfouten}

\begin{itemize}

  \item \textbf{``Het elektrische veld hangt af van de testlading.''}

  Nee. De bronladingen bepalen het veld. De testlading bepaalt welke
  kracht ze in dat veld ondervindt.

  \item \textbf{``Een negatieve testlading keert het veld om.''}

  Nee. Ze keert de zin van de kracht ten opzichte van het veld om.

  \item \textbf{``Veldlijnen zijn echte draden in de ruimte.''}

  Nee. Het zijn voorstellingsmiddelen om een vectorveld zichtbaar te
  maken.

  \item \textbf{``Waar geen veldlijn getekend is, is het veld nul.''}

  Nee. Een veldlijnendiagram toont slechts een selectie van lijnen.

  \item \textbf{``Potentiaal en potentiële energie zijn hetzelfde.''}

  Nee:

  \[
  E_{\mathrm{pot,e}}=qV.
  \]

  \item \textbf{``Als $V=0$, dan is ook $E=0$.''}

  Niet noodzakelijk.

  \item \textbf{``Als $E=0$, dan moet $V=0$.''}

  Ook dat is niet noodzakelijk.

  \item \textbf{``Een positieve en negatieve lading bewegen spontaan
  allebei naar lagere potentiaal.''}

  Nee. Een systeem evolueert spontaan naar lagere potentiële energie.
  Omdat $E_{\mathrm{pot,e}}=qV$, hangt de relatie met $V$ af van het
  teken van $q$.

  \item \textbf{``In een homogeen veld is de potentiaal overal
  dezelfde.''}

  Nee. Het veld is constant, maar de potentiaal verandert lineair in de
  richting van het veld.

\end{itemize}


% ============================================================
% 50. SAMENVATTING
% ============================================================

\FVDsection{Samenvatting}


\FVDsubsection{Elektrische veldsterkte}

\[
\boxed{
\vec E
=
\frac{\vec F_e}{q_t}
}
\]

voor een positieve testlading.

Eenheid:

\[
\boxed{
[E]=\mathrm{\frac NC}.
}
\]


\FVDsubsection{Puntlading}

\[
\boxed{
E
=
k\frac{|Q|}{r^2}.
}
\]

Voor $Q>0$ wijst $\vec E$ van de lading weg.

Voor $Q<0$ wijst $\vec E$ naar de lading toe.


\FVDsubsection{Superpositie}

\[
\boxed{
\vec E_\mathrm{res}
=
\sum_i\vec E_i.
}
\]


\FVDsubsection{Kracht in een veld}

\[
\boxed{
\vec F_e=q\vec E.
}
\]

Voor $q>0$ hebben $\vec F_e$ en $\vec E$ dezelfde zin.

Voor $q<0$ hebben ze tegengestelde zin.


\FVDsubsection{Elektrische potentiële energie}

Voor twee puntladingen:

\[
\boxed{
E_{\mathrm{pot,e}}
=
k\frac{Qq}{r}.
}
\]


\FVDsubsection{Elektrische potentiaal}

\[
\boxed{
V
=
\frac{E_{\mathrm{pot,e}}}{q}.
}
\]

Voor een puntlading:

\[
\boxed{
V
=
k\frac{Q}{r}.
}
\]

Eenheid:

\[
\boxed{
1\,\mathrm V
=
1\,\mathrm{\frac JC}.
}
\]


\FVDsubsection{Superpositie van potentiaal}

\[
\boxed{
V_\mathrm{tot}
=
\sum_iV_i.
}
\]


\FVDsubsection{Spanning}

\[
\boxed{
\Delta V=V_B-V_A.
}
\]

En:

\[
\boxed{
\Delta E_{\mathrm{pot,e}}
=
q\Delta V.
}
\]


\FVDsubsection{Homogeen elektrisch veld}

Voor een verplaatsing evenwijdig aan het veld:

\[
\boxed{
|\Delta V|=Ed.
}
\]

Dus:

\[
\boxed{
E=\frac{|\Delta V|}{d}.
}
\]

Daaruit:

\[
\boxed{
1\,\mathrm{\frac NC}
=
1\,\mathrm{\frac Vm}.
}
\]


\FVDsubsection{Versnelling van een geladen deeltje}

\[
\boxed{
\vec a
=
\frac{q}{m}\vec E.
}
\]


\FVDsubsection{Energieomzetting}

Wanneer alleen de elektrische kracht arbeid verricht:

\[
\boxed{
\Delta E_k
=
-\Delta E_{\mathrm{pot,e}}.
}
\]

Dus:

\[
\boxed{
\Delta E_k
=
-q\Delta V.
}
\]


% ============================================================
% 51. WAT MOET JE ZEKER BEHEERSEN?
% ============================================================

\FVDsection{Wat moet je zeker beheersen?}

Om dit hoofdstuk op het vereiste niveau te beheersen, moet je meer
kunnen dan formules invullen.

Je moet zeker kunnen:

\begin{itemize}

  \item uitleggen wat een elektrisch veld voorstelt;

  \item bronlading en testlading onderscheiden;

  \item de definitie

        \[
        \vec E=\frac{\vec F_e}{q_t}
        \]

        interpreteren;

  \item het veld van een puntlading berekenen met

        \[
        E=k\frac{|Q|}{r^2};
        \]

  \item richting en zin van $\vec E$ bepalen;

  \item veldlijnen tekenen en interpreteren;

  \item uit veldlijndichtheid conclusies trekken over de grootte van
        $E$;

  \item radiale, homogene en dipoolvelden herkennen en vergelijken;

  \item meerdere elektrische velden vectorieel combineren;

  \item symmetrie gebruiken om een elektrisch veld te analyseren;

  \item de kracht op een positieve en negatieve lading bepalen met

        \[
        \vec F_e=q\vec E;
        \]

  \item elektrische potentiële energie interpreteren;

  \item potentiaal berekenen met

        \[
        V=k\frac{Q}{r};
        \]

  \item het verschil tussen $\vec E$, $V$ en
        $E_{\mathrm{pot,e}}$ uitleggen;

  \item de potentiaal van meerdere bronladingen algebraïsch optellen;

  \item uitleggen waarom $V=0$ niet noodzakelijk betekent dat
        $\vec E=0$;

  \item spanning als potentiaalverschil interpreteren;

  \item werken met

        \[
        \Delta E_{\mathrm{pot,e}}=q\Delta V;
        \]

  \item in een homogeen veld het verband

        \[
        |\Delta V|=Ed
        \]

        toepassen;

  \item het verband tussen veldlijnen en equipotentialen uitleggen;

  \item de beweging en energie van een geladen deeltje in een homogeen
        elektrisch veld analyseren;

  \item je berekeningen controleren op eenheden, tekens, richting,
        grootteorde en fysische betekenis.

\end{itemize}

De oefeningen met de essentieel-markering vormen samen de minimale
oefenset.

Let op: ook verschillende oefeningen uit de categorie
\oefU zijn essentieel. Dat is noodzakelijk omdat je niet alleen
standaardtechnieken moet kunnen uitvoeren, maar elektrische velden en
potentialen ook in nieuwe situaties moet kunnen analyseren.


% ============================================================
% 52. SLOT
% ============================================================

\FVDsection{Van elektrische naar magnetische velden}

We hebben elektrische wisselwerkingen nu op verschillende niveaus
leren beschrijven.

Eerst als kracht tussen ladingen:

\[
Q_1,Q_2
\longrightarrow
\vec F_e.
\]

Daarna als elektrisch veld:

\[
Q
\longrightarrow
\vec E
\longrightarrow
\vec F_e.
\]

En ten slotte vanuit energie:

\[
V
\longrightarrow
E_{\mathrm{pot,e}}
\longrightarrow
E_k.
\]

Een belangrijke vraag blijft echter open.

Tot nu toe hebben we vooral elektrische ladingen beschouwd die als
bron van een elektrisch veld optreden.

Maar wat gebeurt er wanneer elektrische ladingen
\textbf{bewegen}?

Een bewegende elektrische lading blijkt nog een ander soort veld te
kunnen veroorzaken:

\[
\boxed{
\text{een magnetisch veld}.
}
\]

Daarmee ontstaat een nieuwe verbinding:

\[
\boxed{
\text{bewegende elektrische lading}
\longrightarrow
\text{magnetisme}.
}
\]

In het volgende thema onderzoeken we hoe elektriciteit en magnetisme
met elkaar verbonden zijn.

\end{document}